% Compiler: pdfLaTeX
\documentclass[11pt,a4paper]{article}

\usepackage[utf8]{inputenc}
\usepackage[T1]{fontenc}
\usepackage{lmodern}
\usepackage{microtype}
\usepackage{geometry}
\usepackage{enumitem}
\usepackage{titlesec}
\usepackage{fancyhdr}
\usepackage{hyperref}

% Document details
\newcommand{\researchtitle}{Sunk cost fallacy}
\newcommand{\researchshorttitle}{Sunk cost fallacy}
\newcommand{\researchdate}{2026-10-25}

% Layout
\geometry{left=2.7cm,right=2.7cm,top=2.5cm,bottom=2.5cm,headheight=14pt}
\hypersetup{hidelinks,pdftitle={\researchtitle}}
\urlstyle{same}
\frenchspacing
\setlength{\parindent}{0pt}
\setlength{\parskip}{0.55em}
\linespread{1.04}
\setlength{\emergencystretch}{1em}

\titleformat{\section}
    {\large\bfseries}{\thesection}{0.65em}{}
\titleformat{\subsection}
    {\normalsize\bfseries}{\thesubsection}{0.65em}{}
\titlespacing*{\section}{0pt}{1.15em}{0.35em}
\titlespacing*{\subsection}{0pt}{0.7em}{0.2em}
\setlist[itemize]{leftmargin=1.35em,itemsep=0.25em,topsep=0.2em}
\setlist[enumerate]{leftmargin=1.6em,itemsep=0.25em,topsep=0.2em}

\pagestyle{fancy}
\fancyhf{}
\fancyhead[L]{\small Random Research}
\fancyhead[R]{\small\researchshorttitle}
\fancyfoot[C]{\small\thepage}
\renewcommand{\headrulewidth}{0pt}
\renewcommand{\footrulewidth}{0pt}

\begin{document}

{\LARGE\bfseries\researchtitle\par}
\vspace{0.4em}
{\small\researchdate\par}

\section{Overview}

A sunk cost is something already spent that cannot be recovered. The sunk
cost fallacy occurs when that past investment makes an inferior option
harder to leave. Money is the obvious example, but time and effort can have
the same pull. \cite{arkes}

The awkward part is that continuing can feel responsible. Walking away
makes the loss visible, even when staying would add another loss.

\section{Findings}

Arkes and Blumer found that people became more willing to continue an
activity after investing in it. In their theatre study, customers who paid
more for season tickets attended more performances during the following
six months. Their experiments also supported the idea that people dislike
appearing wasteful. \cite{arkes}

The investment does not have to be our own. Olivola's experiments found
that people changed their choices in response to costs paid by someone
else. This complicates the familiar advice to ignore what you personally
spent: a gift or somebody else's effort can create a similar sense of
obligation. \cite{olivola}

These findings describe a tendency. A previous investment alone does not
tell us whether continuing a particular project is sensible.

\section{Example}

Suppose someone pays for a course, finishes several lessons, and realizes
that it does not teach what they need. They keep going because the fee
was expensive. If the payment is non-refundable, finishing cannot recover
it. The remaining lessons still cost time.

There may be good reasons to continue, such as a useful qualification or
material that appears later. Those are benefits still available. This
hypothetical example depends on what the course can offer from today.

\section{Discussion}

I find the reluctance to admit waste more interesting than the money.
Stopping can feel like a judgment on the person who started. That makes
a practical decision unnecessarily personal.

Still, I would be careful with advice to quit whenever something becomes
unpleasant. Learning often includes dull periods. My question would be
whether I can explain the expected benefit of continuing without relying
on how much I have already sacrificed.

\section{Practical implications}

For a small decision, I would compare the remaining cost with what I
expect to gain. A useful prompt is: if I had this option available today,
with the next steps unchanged, would I choose it?

Refunds, obligations, and the cost of switching also belong in that
comparison. Calling a cost sunk only makes sense once it really is
unrecoverable.

\renewcommand{\refname}{References}
\begin{thebibliography}{9}
\small
\setlength{\itemsep}{0.35em}

\bibitem{arkes}
Arkes, H. R., and Blumer, C. (1985).
\href{https://www.sciencedirect.com/science/article/pii/0749597885900494}
{The psychology of sunk cost.}
\textit{Organizational Behavior and Human Decision Processes}, 35(1), 124--140.

\bibitem{olivola}
Olivola, C. Y. (2018).
\href{https://doi.org/10.1177/0956797617752641}
{The interpersonal sunk-cost effect.}
\textit{Psychological Science}, 29(7), 1072--1083.

\end{thebibliography}
\end{document}
